\chapter{Sviluppo del firmware per la gestione dell'incubatrice neonatale}
\section{Requisiti del firmware}
Il firmware per il microcontrollore deve essere in grado di gestire in tempo reale ed in maniera concorrente tutti i
componenti del prototipo di incubatrice neonatale, come il pannello informativo frontale, la gestione degli attuatori,
il sensore di temperatura e umidità, la comunicazione seriale con il PC e i controllori per la gestione della temperatura.

Per garantire la corretta sincronizzazione tra i vari processi, non è possibile impiegare codice bloccante come loop
in attesa di soddisfare una certa condizione o l'utilizzo dela funzione \verb|delay()|. Serve usare un approccio basato
su macchine a stati e condizioni che sfruttano la funzione \verb|millis()|. In questo modo, il microcontrollore può
portare avanti più processi in parallelo senza rimanere occupato a concludere un determinato compito e trascurare gli altri.

Siccome il firmware deve operare su un microcontrollore con risorse limitate, l'implementazione deve essere efficiente
sia in termini di complessità computazionale che di memoria. Per questo motivo, si cerca di limitare il più possibile
l'uso di variabili locali e di allocazioni dinamiche nello stack/heap, preferendo l'uso di variabili globali. In questo
modo la memoria RAM utilizzata per le variabili è già conosciuta a tempo di compilazione e risulta facile accorgersi di
possibili problemi di overflow della memoria. Inoltre, l'uso di variabili globali permette di condividere facilmente
lo stato del sistema tra i vari moduli del firmware, senza dover ricorrere continuamente a passaggi di parametri tra
le funzioni.

Infine è richiesto che il codice sia modulare, leggibile, documentato e facilmente modificabile. Nello sviluppo di un
prototipo orientato alla ricerca, infatti, capiterà sicuramente e molto spesso la necessità di modificare il codice
per aggiungere nuove funzionalità, correggere scelte che si rivelano non ottimali o per adattarlo a nuove esigenze.

\section{Configurazione generale}
\begin{itemize}
	\item variabili globali di stato
	\item files e moduli del firmware
	\item istruzioni di compilazione da arduino ide e platform.io
\end{itemize}

\section{Gestione degli attuatori}
\begin{itemize}
	\item controllo manuale o automatico degli attuatori
	\item gestione centralizzata degli attuatori per aggiornare le variabili di stato
	\item gestione della ventola di omogeneizzazione della temperatura
	\item gestione del refill dell'acqua per il sistema di umidificazione
\end{itemize}

\section{Gestione del pannello frontale}
\begin{itemize}
	\item lettura degli switch, con impostazione controllo manuale
	\item lettura dei potenziometri, con filtri per ridurre il rumore
	\item gestione del display lcd con temporizzazioni
	\item gestione degli errori e gestione del led di allarme
\end{itemize}

\section{Gestione della comunicazione seriale al PC}
\begin{itemize}
	\item comunicazioni seriali
	\item codice matlab
\end{itemize}

\section{Gestione del sensore di temperatura e umidità}
\begin{itemize}
	\item gestione non bloccante
	\item utilizzo di una macchina a stati
\end{itemize}
